\documentclass[12pt]{article}

\usepackage{amsmath}
\usepackage{geometry}
\usepackage{hyperref}
\usepackage{parskip}
\usepackage{booktabs}

\geometry{a4paper, margin=1in}

\title{A Latent Log-Ratio Model for Practical Shooting Ratings}
\author{Jay Slater}
\date{\today}

% === Clean shortcuts for Latent Log-Ratio Model ===
\newcommand{\sigmai}{\sigma_i^2}
\newcommand{\ssport}{\sigma_{\mathrm{sport}}^2}      % sport variance
\newcommand{\vstart}{V_{\mathrm{start}}}
\newcommand{\vmax}{V_{\mathrm{max}}}
\newcommand{\sdrift}{\sigma_{\mathrm{drift}}^2}
\newcommand{\psii}{\psi^2}
\newcommand{\xii}{\xi^2}
\newcommand{\omegai}{\omega^2}
\newcommand{\zetai}{\zeta^2}                    % lump coefficient (s^2)
\newcommand{\slump}{\sigma_{\mathrm{lump}}^2}  % event-level lump penalty
\newcommand{\spred}{\sigma_{\mathrm{pred}}^2}
\newcommand{\tref}{T_{\mathrm{ref}}}
\newcommand{\tfloor}{T_{\mathrm{floor}}}

\newcommand{\betad}{\beta}      % dispersion adaptation rate
\newcommand{\gammad}{\gamma}    % surprise adaptation rate
\newcommand{\lambdad}{\lambda}  % momentum adaptation rate
\newcommand{\alphad}{\alpha}    % pairwise blend weight
\newcommand{\rhoc}{\rho}        % cohort intraclass correlation
\newcommand{\nudof}{\nu}       % Student-t degrees of freedom
\newcommand{\ctcut}{c_t}       % Student-t full-strength cutoff
\newcommand{\cthresh}{c}        % baseline robustness threshold
\newcommand{\szero}{s_0}
\newcommand{\kmax}{k_{\mathrm{max}}}
\newcommand{\kappap}{\kappa}    % prediction behavioral scale
\newcommand{\lambdar}{\lambda_{\mathrm{rust}}}
\newcommand{\tgrace}{t_{\mathrm{grace}}}

\newcommand{\xwin}{X_w}
\newcommand{\tmin}{T_{\mathrm{min}}}

\newcommand{\bobs}{B_{\mathrm{obs}}}
\newcommand{\blocal}{B_{\mathrm{local}}}
\newcommand{\bloo}{B_{-i}}
\newcommand{\ephys}{e_i^{\mathrm{phys}}}
\newcommand{\erobust}{e_i^{\mathrm{robust}}}

\DeclareMathOperator*{\median}{Median}


\begin{document}

	\maketitle

	\begin{abstract}
		Traditional rating systems like Elo, Glicko, and TrueSkill were designed for zero-sum or ordinal competitions. Applying them to continuous, proportional-scoring sports (e.g., practical shooting, racing) introduces severe structural friction. Existing systems either discard continuous margins of victory entirely or fall victim to a nuisance-parameter artifact where an outlier winning performance artificially deflates the apparent skill of the entire field.
		
		This paper proposes the Latent Log-Ratio Model, a Bayesian state-space approach that operates natively on continuous, ratio-to-winner scoring. By leveraging the logarithmic structure of multiplicative performances, the model cleanly isolates pairwise skill gaps and exactly cancels the winner's match-day variance.
		
		The theoretical architecture features an idealized $\mathcal{O}(N)$	update per event, explicitly decoupled tracking of epistemic uncertainty and aleatoric performance dispersion, and a globally reversible link function for generating predictions over complete fields. To translate this idealized model to real-world environments, the paper further introduces a robust adaptive filtering framework to handle heavy-tailed outliers, disconnected competitor topologies, and deep-tail scoring non-linearities without introducing structural bias.
	\end{abstract}

	\section{Key Insights and Concepts}
	\label{sec:insights}

	\subsection{The Multiplicative Nature of Scoring}
	In practical shooting, true skill is fundamentally multiplicative, whether measured by hit factor  in USPSA, or total time in time-plus sports like IDPA and ICORE. If Competitor A takes 10 seconds to shoot a stage and Competitor B takes 20 seconds, Competitor A is twice as fast. On a longer stage, the times might be 20 seconds and 40 seconds; the absolute margin changes, but the $2\times$ ratio remains constant.

	\subsection{The ``Winner's Curse'' in Percentage Scoring}
	Matches are scored as percentages of the top competitor's performance. Using these raw percentages injects the winner's match-day variance into the entire field, a phenomenon we term the ``winner's curse'' (distinct from the overbidding phenomenon in auction theory).
	\begin{itemize}
		\item \textbf{Hit Factor} (USPSA): Each competitor's stage score is their points divided by their time (the ``hit factor''). The stage percentage is this hit factor divided by the winner's hit factor. On each stage, competitors receive a fraction of the total points available from that stage determined by their finish ratio in \([0, 1]\). A match score is the sum of all stage scores.
		\item \textbf{Time-Plus} (IDPA/ICORE): Each competitor's stage score is their raw time plus penalties. The implied percentage is the winner's time divided by the competitor's time.
	\end{itemize}
	If the match winner has an exceptionally good day, every other competitor's percentage drops, mathematically suggesting the field performed poorly when they did not.

	\subsection{The Logarithmic Bridge}
	Let $P_i$ denote competitor $i$'s \textit{latent} (true) performance parameter on any positive ratio scale where larger is better. On match day, each competitor produces a realized performance $X_i$---a noisy draw reflecting their current skill. The scoreboard reports only the ratio of realized performances:
	\begin{equation}
		S_i = \frac{X_i}{\xwin}
	\end{equation}
	where $\xwin = \max_j X_j$ is the realized performance of whoever won the match. Critically, neither $P_i$ nor $X_i$ is directly observed; only the ratio $S_i$ is published.

	For any two competitors $A$ and $B$, the difference of their log-percentages is
	\begin{equation}
		\ln(S_A) - \ln(S_B) = \ln\!\left(\frac{X_A}{\xwin}\right) - \ln\!\left(\frac{X_B}{\xwin}\right)
	\end{equation}
	Using $\ln(x/y) = \ln(x) - \ln(y)$, this becomes
	\begin{equation}
		\ln(S_A) - \ln(S_B) = (\ln(X_A) - \ln(\xwin)) - (\ln(X_B) - \ln(\xwin))
	\end{equation}
	and the common anchor cancels exactly:
	\begin{equation}
		\ln(S_A) - \ln(S_B) = \ln(X_A) - \ln(X_B) = \ln\!\left(\frac{X_A}{X_B}\right)
	\end{equation}
	This is the central trick of the model: the realized winning performance $\xwin$ appears only as a nuisance normalization constant in the published percentages; it does not survive pairwise comparison in log space, an idea derived from compositional data analysis\cite{aitchison1982compositional}\footnote{Although this model leverages log-ratio transformations that are fundamental to compositional data analysis, match percentages do not themselves constitute compositional data by the strict definition, as they violate the unit-sum constraint. The data do, however, have scale invariance: the absolute magnitudes are irrelevant or (in our case) misleading, and all the usable information is contained strictly in the \emph{ratios} between the components.}. The resulting quantity is the pure ratio of the competitors' realized performances, independent of who won the match or by how much.

	Because each realized performance is a noisy draw from around the competitor's latent skill, the log-ratio decomposes as:
	\begin{equation}
		\ln\!\left(\frac{X_A}{X_B}\right) = \ln\!\left(\frac{P_A}{P_B}\right) + (\varepsilon_A - \varepsilon_B)
	\end{equation}
	where $\varepsilon_i$ represents competitor $i$'s match-day noise. In expectation, the noise cancels and the observed log-ratio is a direct estimate of the latent skill gap:
	\begin{equation}
		E\!\left[\ln\!\left(\frac{X_A}{X_B}\right)\right] = \ln\!\left(\frac{P_A}{P_B}\right) \approx R_A - R_B
	\end{equation}
	In time-plus sports, this same derivation applies by defining $X_i$ as a speed $X_i = \frac{1}{T_i}$, which yields $S_i = \frac{X_i}{\xwin} = \frac{1/T_i}{1/\tmin} = \frac{\tmin}{T_i}$ and $\ln(S_A) - \ln(S_B) = \ln\left(\frac{1/T_A}{1/T_B}\right) = \ln\left(\frac{T_B}{T_A}\right)$.

	\subsection{Ratings as Expected Log-Ratios}
	This cancellation property leads directly to the model's core identity. Since ratings $R_i$ track each competitor's current latent log-performance $\ln(P_i)$, the difference between any two ratings is the expected log-ratio of their realized performances:
	\begin{equation}
		R_i - R_j = E[\ln(S_i) - \ln(S_j)] = E\!\left[\ln\!\left(\frac{X_i}{X_j}\right)\right] = \ln\!\left(\frac{P_i}{P_j}\right)
	\end{equation}

	Equivalently, the median ratio of two competitors' realized performances is:
	\begin{equation}
		\median\!\left(\frac{X_i}{X_j}\right) = \exp(R_i - R_j)
	\end{equation}

	The rating difference \textit{is} the expected difference in log-performance: a rating gap of $\ln(2) \approx 0.69$ means one competitor's latent skill is twice the level of the other. Deep-tail pathologies are handled by inflating observation noise for very weak finishes, not by deforming the log-ratio mean link, so the prediction mapping remains globally reversible through the exponential function.
	
	\section{The Core State-Space Architecture}
	\label{sec:core}
	
	\subsection{Generative Model}
	The Latent Log-Ratio model formulates competitor skill as a Bayesian state-space tracking problem on a logarithmic manifold. Each competitor $i$ possesses a true, time-varying latent performance parameter $P_i \in \mathbf{R}^+$. On match day, the competitor produces a realized performance $X_i$ whose logarithm is a noisy realization centered on their latent skill:
	\begin{equation}
		\ln(X_i) = \ln(P_i) + \varepsilon_i, \qquad \varepsilon_i \sim \mathcal{N}(0,\, \ssport + \sigmai)
	\end{equation}
	The stochastic perturbation $\varepsilon_i$ combines the irreducible, environmental noise of the sport ($\ssport$) with competitor $i$'s individual performance dispersion ($\sigmai$). 
	
	The published scoreboard does not reveal $X_i$ or $P_i$; it reports only the normalized ratio relative to the event winner:
	\begin{equation}
		S_i = \frac{X_i}{\xwin}, \qquad \xwin = \max_{j} X_j
	\end{equation}
	The rating system maintains a Gaussian posterior over each competitor's latent log-performance $\ln(P_i)$, parameterized by a mean rating $R_i$ and an epistemic variance $V_i$.
	
	\subsection{Competitor State Representation}
	Each competitor is parameterized by a four-tuple state $(R_i, V_i, \sigmai, m_i)$:
	\begin{itemize}
		\item \textbf{$R_i$ (Rating):} The posterior mean estimate of $\ln(P_i)$. An arbitrary baseline competitor can be anchored at $R_0 = 0$.
		\item \textbf{$V_i$ (Epistemic Variance):} The variance of the posterior rating distribution, quantifying system uncertainty regarding $R_i$.
		\item \textbf{$\sigmai$ (Performance Dispersion):} An estimate of the competitor's aleatoric variability—the degree to which their match-day performance naturally fluctuates around their latent mean.
		\item \textbf{$m_i$ (Innovation Momentum):} An exponential moving average of recent prediction residuals, tracking persistent directional skill evolution (breakouts or declines).
	\end{itemize}
	
	\subsection{The Idealized Update Cycle}
	The core update loop operates under the assumption of independent Gaussian perturbations, complete graph connectivity, and well-behaved ratio scoring. Match results are processed through an $\mathcal{O}(N)$ Bayesian filter that decouples the estimation of the match anchor from the state update of each individual competitor.
	
	\subsubsection*{Step 1: Epistemic State Aging}
	Between competitive events, true skill undergoes an unobserved random walk. Prior to evaluating a match, epistemic uncertainty expands linearly with elapsed time $\Delta t$:
	\begin{equation}
		V_i^{\mathrm{age}} = V_i + \sdrift \cdot \Delta t
	\end{equation}
	where $\sdrift$ is the continuous skill drift rate.
	
	\subsubsection*{Step 2: Logarithmic Transformation}
	Match percentages are projected onto the logarithmic difference scale:
	\begin{equation}
		L_i = \ln(S_i) = \ln(X_i) - \ln(\xwin)
	\end{equation}
	Because $S_i \le 1.0$, log-scores are strictly non-positive ($L_i \le 0$).
	
	\subsubsection*{Step 3: Consensus Anchor Estimation ($B_{-i}$)}
	To resolve the unobserved winning performance $\ln(\xwin)$, we evaluate the field's expected performance residuals:
	\begin{align}
		a_i &= R_i - L_i \nonumber \\
		&= R_i - (\ln(X_i) - \ln(\xwin)) \nonumber \\
		&= (R_i - \ln(X_i)) + \ln(\xwin)
	\end{align}
	Because $E[R_i - \ln(X_i)] = 0$, each residual $a_i$ serves as an unbiased, noisy observation of the nuisance parameter $\ln(\xwin)$. 
	
	Competitors are weighted inversely by their total expected variance:
	\begin{equation}
		W_i = \frac{1}{V_i^{\mathrm{age}} + \ssport + \sigmai}
	\end{equation}
	To prevent self-reinforcing feedback during the individual update, we form a leave-one-out precision-weighted consensus estimate of the anchor for each competitor:
	\begin{equation}
		B_{-i} = \frac{\sum_{j \neq i} W_j a_j}{\sum_{j \neq i} W_j}
	\end{equation}
	Under independent prior uncertainties, the structural variance of this consensus baseline is:
	\begin{equation}
		\mathrm{Var}(B_{-i}) = \frac{1}{\sum_{j \neq i} W_j}
	\end{equation}
	
	\subsubsection*{Step 4: Reconstruction of Absolute Performance}
	By adding the leave-one-out consensus anchor to the competitor's log-percentage, the nuisance anchor cancels:
	\begin{equation}
		O_i = L_i + B_{-i} \approx \ln(X_i)
	\end{equation}
	The quantity $O_i$ provides a direct, noisy measurement of the competitor's absolute realized log-performance on the global rating scale. The total observation variance associated with $O_i$ compounds the physical match noise, individual dispersion, and the estimation variance of the anchor:
	\begin{equation}
		R_{\mathrm{obs},i} = \ssport + \sigmai + \mathrm{Var}(B_{-i})
	\end{equation}
	
	\subsubsection*{Step 5: Trend-Adjusted Kalman Mean Update}
	Before updating the rating mean, the competitor's committed innovation momentum $m_i$ is leveraged as an adaptive estimate of unmodeled velocity. We project this momentum into adaptive process noise:
	\begin{equation}
		\text{Trend}_i = \frac{m_i^2}{\lambdad}
	\end{equation}
	where $\lambdad$ is the momentum adaptation rate. This adaptive process noise is injected directly into the prior variance:
	\begin{equation}
		V_i^{\mathrm{prior}} = V_i^{\mathrm{age}} + \gammad \cdot \text{Trend}_i
	\end{equation}
	where $\gammad$ governs the sensitivity of the filter to structural state shifts. 
	
	The filter computes the innovation $e_i$, total innovation variance $T_i$, and the Kalman gain $K_i$:
	\begin{equation}
		e_i = O_i - R_i, \qquad T_i = V_i^{\mathrm{prior}} + R_{\mathrm{obs},i}, \qquad K_i = \frac{V_i^{\mathrm{prior}}}{T_i}
	\end{equation}
	The rating mean is updated via the linear Kalman step:
	\begin{equation}
		R_i \leftarrow R_i + K_i e_i
	\end{equation}
	
	\subsubsection*{Step 6: Shock Adaptation and Epistemic Variance Contraction}
	Following the mean update, the system evaluates whether the observation contained an unmodeled variance shock. We isolate the residual excess variance:
	\begin{equation}
		\text{Shock}_i = \max\!\left(0,\; e_i^2 - T_i\right)
	\end{equation}
	The posterior epistemic variance updates via standard Kalman contraction, augmented by the instantaneous shock:
	\begin{equation}
		V_i \leftarrow V_i^{\mathrm{prior}}(1 - K_i) + \gammad \cdot \text{Shock}_i
	\end{equation}
	This architecture enforces an asymmetric operational split between sustained trends and sudden anomalies:
	\begin{itemize}
		\item \textbf{Trend} is injected into the \textit{prior} variance $V_i^{\mathrm{prior}}$, raising $K_i$ and permitting the rating mean $R_i$ to rapidly track sustained skill breakouts.
		\item \textbf{Shock} is injected into the \textit{posterior} variance $V_i$, widening epistemic uncertainty following a severe surprise without forcing the current mean to overreact to an unverified single-event anomaly.
	\end{itemize}
	
	\subsubsection*{Step 7: Momentum and Dispersion Tracking}
	Finally, the non-Kalman auxiliary states are updated via exponential moving averages. The innovation momentum updates to track persistent directional bias:
	\begin{equation}
		m_i \leftarrow (1 - \lambdad) m_i + \lambdad e_i
	\end{equation}
	To update the aleatoric dispersion $\sigmai$, we isolate behavioral inconsistency by subtracting the directional trend $m_i$ and accounting for expected measurement variances:
	\begin{equation}
		E_i = \max\!\left(0,\; (e_i - m_i)^2 - V_i^{\mathrm{age}} - \ssport\right)
	\end{equation}
	The individual performance dispersion adapts according to:
	\begin{equation}
		\sigmai \leftarrow (1 - \betad)\sigmai + \betad E_i
	\end{equation}
	Subtracting $m_i$ prevents monotonic skill trajectories from being misclassified as behavioral volatility, ensuring that rapid learning curves do not artificially inflate the aleatoric dispersion parameter.
	
	\section{Robust Extensions for Real-World Match Topologies}
	\label{sec:robust}
	
	The idealized state-space model derived in Section 2 operates on three fundamental assumptions: that performance innovations are strictly Gaussian, that the competitor graph is topologically dense and well-connected, and that scoring behaves as a scale-invariant ratio process. 
	
	In production environments, real-world match data routinely violates all three assumptions. Competitors experience catastrophic physical equipment failures, unrated cohorts compete in local isolation, and certain scoring systems inject additive discrete penalties into multiplicative time scales. This section introduces targeted robustifications designed to handle these pathologies without corrupting the underlying log-linear rating link. The corresponding insertions into the update loop are given in Section~\ref{sec:algorithm}.
	
	\subsection{Heavy Tails and Contaminated Innovations}
	\label{sec:heavytails}
	The Gaussian likelihood assumed in Section~2 is thin-tailed. Catastrophic gear failures (jams, squibs, DQs) inject 4--8\,$\sigma$ outliers into both the consensus anchor and the individual Kalman update. A pure $L^2$ estimate of $B_{-i}$ lets a single broken gun drag $\ln(\xwin)$ for the entire field; an undamped Kalman mean chases the same event into $R_i$.
	
	Two M-estimators isolate the mean path from the variance path. First, the consensus anchor is Huber-weighted. A raw, prior-regularized center $B_{\mathrm{raw}} = (\sum_i W_i a_i)/(\sum_i W_i)$ standardizes residuals, and weights taper beyond the robustness threshold $\cthresh$:
	\begin{equation}
		r_i = \frac{a_i - B_{\mathrm{raw}}}{\sqrt{1/W_i}}, \qquad
		\tilde{W}_i = W_i \cdot \min\!\left(1, \frac{\cthresh}{|r_i|}\right).
	\end{equation}
	Leave-one-out $B_{-i}$ and $\mathrm{Var}(B_{-i})$ then use $\tilde{W}_j$ rather than $W_j$.
	
	Second, the rating mean is a Student-$t$ M-estimator. With total innovation uncertainty $T_i = V_i^{\mathrm{prior}} + R_{\mathrm{obs},i}$ and $z_i = e_i / \sqrt{T_i}$,
	\begin{equation}
		w_t(z_i) = \min\!\left(1,\frac{\nudof + \ctcut^2}{\nudof + z_i^2}\right), \qquad
		R_i \leftarrow R_i + K_i \cdot w_t(z_i) \cdot e_i.
	\end{equation}
	Setting $\ctcut = 1$ recovers the standard Student-$t$ kernel; larger values admit a wider full-strength band before tapering. The cutoff may be split by the sign of $e_i$ when contamination is known to be one-sided. Crucially, $w_t$ applies \emph{only} to the rating mean. Shock and the dispersion EMA continue to process physical anomalies under a trustworthy log-link, so a genuine equipment failure can still shatter epistemic confidence without whipping $R_i$ around.
	
	\subsection{Network Sparsity and Topological Isolation}
	\label{sec:sparsity}
	The independent-prior assumption $\rhoc = 0$ treats the field's ratings as rigidly anchored to global truth. In reality, a local club of unrated shooters forms an unanchored floating subgraph. Under independent priors, $\mathrm{Var}(B_{-i}) = 1/\sum W_j$ collapses toward zero as $N$ grows even when every competitor is unknown, producing false epistemic certainty and explosive rating inflation for early winners (the ``Day 1 rating explosion'').
	
	Assuming a uniform prior correlation $\rhoc$ among the cohort's epistemic uncertainties yields an exact $\mathcal{O}(N)$ expansion of baseline variance under equicorrelated priors:
	\begin{equation}
		\mathrm{Var}(B_{-i}) = \frac{1}{\sum_{j \neq i} \tilde{W}_j} + \frac{\rhoc}{\left(\sum_{j \neq i} \tilde{W}_j\right)^2} \left[ \left( \sum_{j \neq i} \tilde{W}_j \sqrt{V_j^{\mathrm{age}}} \right)^2 - \sum_{j \neq i} \tilde{W}_j^2 V_j^{\mathrm{age}} \right].
	\end{equation}
	The second term asymptotes to a nonzero floor for unrated fields and vanishes as competitors mature ($V_j \to 0$). A large field of unknowns becomes a \emph{floating} anchor rather than a rigid datum, suppressing false Kalman confidence.
	
	Even with $\rhoc > 0$, $V_i$ still contracts with repeated observations inside an island: statistical certainty within a local cohort outpaces topological mixing with the global graph. Treat graph isolation as a reliability penalty. Let $k_i$ be competitor $i$'s history length and
	\begin{equation}
		\mu_i = \min\!\left(1.0,\; \frac{\ln(k_i + 1)}{\ln(\kmax + 1)}\right), \qquad
		\bar{\mu} = \frac{\sum_j W_j \mu_j}{\sum_j W_j}.
	\end{equation}
	
	Field novelty is
	\begin{equation}
		\sigma_{\mathrm{novel}}^2 = \psii \cdot (1 - \bar{\mu}).
	\end{equation}
	
	When $\bar{\mu} = 0$, the penalty reaches $\psii$ and throttles update velocity for the entire cohort; as the field matures, it decays to zero.
	
	Global $B_{-i}$ averages field difficulty but discards whom each competitor beat under shared conditions. For an opponent subset $\mathcal{S}_i$ (e.g., the same schedule, or nearby rating/finish), form $B_{\mathrm{local}}$ analogously to $B_{-i}$ and blend
	\begin{equation}
		\bobs = (1 - \alpha_i)\, B_{-i} + \alpha_i\, \blocal, \qquad
		O_i = L_i + \bobs.
	\end{equation}
	On well-mixed large fields, $\blocal$ often agrees with $B_{-i}$ and $\alphad = 0$ recovers a pure $\mathcal{O}(N)$ global update. Nonzero $\alphad$ preserves local peer comparisons across stratified multi-day schedules (e.g., a super squad in rain versus the rest of the field in sun). Blend uncertainty propagates as
	\begin{equation}
		\mathrm{Var}(\bobs) = (1 - \alpha_i^2)\, \mathrm{Var}(B_{-i}) + \alpha_i^2\, \mathrm{Var}(\blocal).
	\end{equation}
	
	\subsection{Scoring Pathologies and Stage Geometry}
	\label{sec:pathologies}
	\subsubsection*{A. Duration-Dependent Heteroskedasticity (Additive Lumps in Time-Plus)}
	\label{sec:lumps}
	The core model assumes that performance is scale-invariant across stage lengths: a competitor who is $10\%$ faster takes $10\%$ less time on an 8-second classifier and a 40-second field course alike. In time-plus disciplines (e.g., ICORE, IDPA), this multiplicative assumption holds well for movement pace and per-shot hit accuracy (where the count of dropped points scales with the number of shots, keeping expected penalty time roughly proportional to duration). It fails, however, for discrete additive blunders.
	
	Fixed-time blunders---a fumbled reload, a procedural penalty, or isolated dropped shots---act as additive \textit{lumps} $P$ measured in seconds. If a competitor produces a total time $T = T_{\mathrm{raw}} + P$, a first-order Taylor expansion gives the log-space contribution
	\begin{equation}
		\ln(T) = \ln\left(T_{\mathrm{raw}}\left(1 + \frac{P}{T_{\mathrm{raw}}}\right)\right) \approx \ln(T_{\mathrm{raw}}) + \frac{P}{T_{\mathrm{raw}}}.
	\end{equation}
	Because the penalty is divided by raw duration, its variance contribution in log-score space scales as
	\begin{equation}
		\mathrm{Var}\!\left(\frac{P}{T_{\mathrm{raw}}}\right) \approx \frac{\mathrm{Var}(P)}{T_{\mathrm{raw}}^2}.
	\end{equation}
	On an 8-second classifier, an isolated two-second blunder introduces a large log-space residual ($P/T \approx 0.25$). If judged against a constant sport noise $\ssport$, this error severely breaches total innovation uncertainty $T_i$. While the Student-$t$ M-estimator $w_t(z_i)$ protects the rating mean $R_i$, unmodeled observation variance causes the filter's single-event $\text{Shock}_i$ to fire, misclassifying an aleatoric execution slip as an epistemic collapse. On the subsequent event, the competitor enters with an inflated prior variance and an elevated Kalman gain, so ordinary noise can yank the rating mean. This pathology is not unique to stage-level tracking: short matches evaluated at the aggregate match level, such as three-stage classifier specials, can suffer the same $1/T_{\mathrm{match}}^2$ magnification, though empirically, the likelihood of isolated errors of this nature drops in proportion to the length of a given match or stage.
	
	To insulate the filter, we introduce a duration-dependent lump penalty with compact support. Rather than conditioning on individual times $T_i$---which would endow slower competitors with spuriously quieter observations---we treat the winning time of the rated event (stage or match) $\tmin$ as an exogenous instrument for course length. Let $T_{\mathrm{eff}} = \max(\tmin,\, \epsilon)$, where $\epsilon \approx 0.5$--$1.0\,\mathrm{s}$ is a numerical floor so that $1/T_{\mathrm{eff}}^2$ remains defined when $\tmin$ is at or near zero; non-positive times themselves are handled by the padding map in D. Using a reference duration $\tref$ (typically $20$--$30\,\mathrm{s}$ for ICORE) representing a course length where additive blunders are sufficiently diluted by movement time, the lump variance penalty is
	\begin{equation}
		\slump(T) = \zetai \cdot \max\!\left(0,\; \frac{1}{T_{\mathrm{eff}}^2} - \frac{1}{\tref^2}\right),
	\end{equation}
	where $\zetai = \mathrm{Var}(P)$ is the variance of the additive lump \emph{in clock seconds}, not the log-space quantity $\mathrm{Var}(P/T_{\mathrm{raw}})$. $\slump(T)$ is the \emph{excess} lump noise of a short event relative to a medium field course. In implementation terms, $\zeta$ is the RMS size of the blunders this term is meant to absorb: a characteristic procedural or fumble of $2\,\mathrm{s}$ suggests a starting $\zetai \approx 4\,\mathrm{s}^2$, scaled down if such lumps are infrequent. For events where $\tmin \ge \tref$, the term evaluates to identically zero, preserving the calibrated sport variance $\ssport$ as the exact floor on standard field courses. The same $Q_i$ gate necessarily throttles genuine short-event pace signal (for example a slow draw on a speed shoot); that is the intended information-content tradeoff, not a deformation of the mean link.
	
	Ontologically, $\slump(T)$ belongs strictly to the \textit{reliability / misspecification} category, entering the system identically to deep-tail noise $\eta(S_i)$ and weak-field pathology $\sigma_{\mathrm{path}}^2$:
	\begin{enumerate}
		\item \textbf{Compression of baseline weights.} The lump penalty is injected into the denominator of the competitor weights:
		\begin{equation}
			W_i = \frac{1}{V_i^{\mathrm{age}} + \ssport + \sigmai + \eta(S_i) + \slump(T)}.
		\end{equation}
		On short events, $\slump(T)$ compresses the dynamic range of precisions. This prevents a single competitor's clean or fortunate run on an 8-second classifier from exerting disproportionate leverage over the consensus match anchor $B_{-i}$.
		
		\item \textbf{Joint Shock and dispersion suppression via $Q_i$.} The term is added to the total observation variance $R_{\mathrm{obs},i}$ \textit{outside} the clean physical variance $R_{\mathrm{clean},i}$:
		\begin{equation}
			R_{\mathrm{obs},i} = R_{\mathrm{clean},i} + \sigma_{\mathrm{novel}}^2 + \eta(S_i) + \sigma_{\mathrm{path}}^2 + \slump(T).
		\end{equation}
		This placement operates simultaneously on both sides of the shock detector. While $T_i = V_i^{\mathrm{prior}} + R_{\mathrm{obs},i}$ expands, the Observation Quality factor $Q_i = R_{\mathrm{clean},i} / R_{\mathrm{obs},i}$ contracts, attenuating the physical innovation $\ephys = e_i \cdot Q_i$. In the shock evaluation
		\begin{equation}
			\text{Shock}_i = \max\!\left(0,\; (\ephys)^2 - T_i\right),
		\end{equation}
		$(\ephys)^2$ is suppressed as $T_i$ rises, so routine short-event blunders do not detonate epistemic variance. Because the dispersion driving residual $E_i$ also relies on $\ephys$, the quality gate prevents those lumps from permanently corrupting $\sigmai$.
	\end{enumerate}
	
	\subsubsection*{B. Deep-Tail Noise ($\eta(S_i)$) and Score Flooring}
	As $S_i \to 0$, log-scores become geometrically extreme ($\ln 0.1 = -2.3$, $\ln 0.01 = -4.6$). Very weak finishes typically reflect catastrophic breakdown---penalties, stage failures, or far-tail geometry---not proportional skill. Compressing $L_i$ itself would destroy the model's log-linear predictive semantics. Instead the mean link stays intact and the observation is treated as an imprecise instrument.
	
	Floor percentages to keep the logarithm defined, then apply a piecewise quadratic reliability penalty below the threshold $\szero$:
	\begin{equation}
		L_i = \ln(\max(S_{\mathrm{floor}}, S_i)), \qquad
		\eta(S_i) =
		\begin{cases}
			0, & S_i \ge \szero \\
			\xii \left(\frac{\szero - S_i}{\szero}\right)^2, & S_i < \szero.
		\end{cases}
	\end{equation}
	The term $\eta(S_i)$ is not generative $\varepsilon_i$; it is score-dependent observation-model skepticism. Dual roles:
	\begin{enumerate}
		\item In weights: $W_i \leftarrow 1/(V_i^{\mathrm{age}} + \ssport + \sigmai + \eta(S_i) + \slump(T))$, so deep tails cannot dominate the consensus anchor.
		\item In own observation variance: $R_{\mathrm{obs},i} \leftarrow R_{\mathrm{obs},i} + \eta(S_i)$, hence in the $Q_i$ denominator.
	\end{enumerate}
	Finishes above $\szero$ are unaffected. Large mid-pack residuals under a trustworthy log-link remain fully available to Shock.
	
	\subsubsection*{C. Small and Bottom-Heavy Fields ($\sigma_{\mathrm{path}}^2$)}
	\label{sec:weakfield-overview}
	Tiny matches in which many non-winners finish far behind the winner are not precise instruments for elite skill separation. Define a size taper that decays to zero at $n_{\mathrm{max}}$,
	\begin{equation}
		f_{\mathrm{size}}(N) = \max\!\left(0,\frac{n_{\mathrm{max}} - N}{n_{\mathrm{max}} - 2}\right),
	\end{equation}
	let $\mathcal{W}$ be the set of non-winners whose scores fall below $s_w$, and write $q = |\mathcal{W}|/(N-1)$ for the weak fraction. If $q < q_0$, no penalty applies. Otherwise a depth severity
	\begin{equation}
		g = \min\!\left(1,\frac{1}{|\mathcal{W}| \cdot (-\ln s_w)} \sum_{j \in \mathcal{W}} \max\!\left(0,\ln\!\left(\frac{s_w}{\max(S_{\mathrm{floor}}, S_j)}\right)\right)\right)
	\end{equation}
	combines with the size taper into a match-level reliability penalty of the same ontological category as novelty and deep-tail noise:
	\begin{equation}
		\sigma_{\mathrm{path}}^2 = \omegai \cdot \bigl(f_{\mathrm{size}}(N) \cdot q \cdot g\bigr).
	\end{equation}
	This term enters $R_{\mathrm{obs},i}$ for the entire field (and therefore reduces $Q_i$): ``this match may be real, but it is not a precise instrument.''
	
	\subsubsection*{D. Singularity Hygiene: Time Padding for Bonus Rules}
	\label{sec:padding}
	Lump variance (A) handles ordinary short events by inflating observation noise. It does not define $\ln$ of a non-positive time. Rule sets that subtract fixed time for accuracy (ICORE bonus rings) can produce near-zero or negative stage times, converting the data from a ratio scale to an interval scale.
	
	When the winning time $\tmin < \tfloor$, add the translation $k = \tfloor - \tmin$ to every competitor's time before forming percentages:
	\begin{equation}
		T_i' = T_i + k.
	\end{equation}
	All $T_i'$ then satisfy $T_i' \ge \tfloor > 0$, rank order is preserved, and the logarithm remains defined. Padding is a mean-link singularity patch, not a substitute for $\slump(T)$. Universal rescaling of long field courses down to a reference time is not used: it would impose a ``speed tax'' on proportional pace (a genuine $10\%$ deficit becomes a $25\%$ ratio penalty on a $50\,\mathrm{s}$ stage mapped to a $20\,\mathrm{s}$ reference) and would discard the physical duration of the task.
	
	\subsection{Information Gating and Filter Safeguards}
	\label{sec:gating}
	Reliability penalties (novelty, tail noise, weak-field pathology, lump variance) must throttle the Kalman mean without corrupting the auxiliary states that have no built-in gain denominator. Partition observation variance into generative physical noise---quantities that belong in the match-day likelihood under the log-ratio model---and reliability / misspecification penalties that express skepticism about whether the observation's magnitude is a precise instrument for latent skill:
	\begin{equation}
		R_{\mathrm{clean},i} = \ssport + \tilde{\sigma}_i^2 + \mathrm{Var}(\bobs),
	\end{equation}
	\begin{equation}
		R_{\mathrm{obs},i} = R_{\mathrm{clean},i} + \sigma_{\mathrm{novel}}^2 + \eta(S_i) + \sigma_{\mathrm{path}}^2 + \slump(T),
	\end{equation}
	\begin{equation}
		Q_i = \frac{R_{\mathrm{clean},i}}{R_{\mathrm{obs},i}} \in (0,1].
	\end{equation}
	Kalman gain $K_i = V_i^{\mathrm{prior}} / (V_i^{\mathrm{prior}} + R_{\mathrm{obs},i})$ discounts \emph{all} observation noise. $Q_i$ strips only the reliability component before feeding momentum, dispersion, and Shock:
	\begin{equation}
		\ephys = e_i \cdot Q_i, \qquad \erobust = e_i \cdot w_t(z_i) \cdot Q_i.
	\end{equation}
	Large residuals under a trustworthy log-link still reach the shock detector; unanchored fields, geometrically unreliable deep tails, and duration-misspecified short events do not permanently inflate $m_i$ or $\sigmai$.
	
	Aleatoric throttling is subordinated to epistemic certainty so that a highly uncertain competitor is not shielded from updates by their own dispersion tracker:
	\begin{equation}
		c_i = 1 - \min\!\left(1,\frac{V_i^{\mathrm{age}}}{\vmax}\right), \qquad
		\tilde{\sigma}_i^2 = \sigmai \cdot c_i,
	\end{equation}
	\begin{equation}
		\lambdad_i^{\mathrm{eff}} = \lambdad \cdot \max(0.25, c_i), \qquad
		\betad_i^{\mathrm{eff}} = \betad \cdot \max(0.25, c_i).
	\end{equation}
	Veterans contribute full behavioral dispersion to $R_{\mathrm{obs}}$; unrated or long-inactive competitors remain fluid. The dispersion EMA target is itself clamped to prevent aleatoric explosions from a single residual. Let $E_i$ be the trend-adjusted squared residual of Section~2; then
	\begin{equation}
		\sigma_{i,\mathrm{ref}}^2 = 0.25\,\ssport, \qquad
		q_i = \mathrm{clamp}\!\bigl(E_i,\; 0.001\,\vstart,\; 9\cdot\max(\sigmai,\, \sigma_{i,\mathrm{ref}}^2)\bigr),
	\end{equation}
	so even a physical outlier under a trustworthy log-link cannot immediately armor the competitor in $\sigmai$.
	
	\subsection{Historical Non-Stationarity (Skill Rust)}
	\label{sec:rust}
	A strict forward-only filter freezes early pioneers against the initial competitor pool. As the sport matures, the filter resolves that early demographic downward; retirees who left before the datum shifted appear mathematically superior to active modern champions. Logistic systems hide the artifact by saturating margins of victory; continuous log-ratios expose it. Rather than a backward smoothing pass, expected ratings revert multiplicatively toward the baseline $R_{\mathrm{start}} = 0$ after a grace period that absorbs ordinary off-season gaps:
	\begin{equation}
		\Delta t_{\mathrm{eff}} = \max(0,\, \Delta t - \tgrace), \qquad
		R_i^{\mathrm{age}} = R_{\mathrm{start}} + (R_i - R_{\mathrm{start}}) \cdot \exp(-\lambdar \cdot \Delta t_{\mathrm{eff}}).
	\end{equation}
	A competitor at the baseline experiences zero drift; an elite outlier slowly regresses. The same map also models perishable fine-motor skill. In the update loop, $R_i^{\mathrm{age}}$ replaces $R_i$ in baseline residuals, pairwise expected margins, the innovation, the Kalman mean update, and out-of-sample predictions.

	\section{System Architecture and Algorithm}
	\label{sec:algorithm}

	This section specifies the production filter in execution order. The generative identity, winner-nuisance cancellation, and Kalman Trend/Shock split are developed in Section~\ref{sec:core}; the robustness terms that appear in-line below are motivated in Section~\ref{sec:robust}. Those derivations are not restated.

	\subsection{Inputs, State, and Initialization}
	Each competitor is the four-tuple $(R_i, V_i, \sigmai, m_i)$ of Section~\ref{sec:core}. Unrated competitors are initialized at $R_i = 0$, $V_i = \vstart$, $m_i = 0$, and $\sigmai = 0.001\,\vstart$. An event supplies a field of finish ratios $S_i \in (0,1]$, the event date, and, for time-plus scoring, the winning time $\tmin$. Competitors with $S_i = 0$ (DNF) are omitted from the update.

	The global baseline is $\mathcal{O}(N)$. A nonempty local opponent set of typical size $k$ raises the cost to $\mathcal{O}(N\cdot k)$.

	Throughout, $R_i^{\mathrm{age}}$ of Step~1 replaces $R_i$ in residuals, innovations, and the Kalman mean; the committed mean is written only at Step~8. Epistemic aging of $V_i$ is likewise applied on entry to the event, not as a post-commit rewrite.

	\subsection{Model Parameters}
	Tables~\ref{tab:params-variance} and \ref{tab:params-dimensionless} list the production knobs. Full ranges and a tuning sequence appear in Appendices~\ref{app:parameters} and \ref{app:tuning}.

	All entries in Table~\ref{tab:params-variance} are log-variance units scaled to the sport-variance anchor $1.0\,\mathrm{SV}$, except \(\zetai\), which is a time-domain variance in $\mathrm{s}^2$. Table~\ref{tab:params-dimensionless} collects rates, weights, and thresholds.

	\begin{table}[htbp]
		\centering
		\caption{Parameters expressed in variance units.}
		\label{tab:params-variance}
		\begin{tabular}{llp{7cm}}
			\toprule
			\textbf{Parameter} & \textbf{Symbol} & \textbf{Role} \\
			\midrule
			Sport variance & \(\ssport\) & Irreducible noise of a single match observation (fundamental anchor, defined as 1.0 SV) \\
			Starting variance & \(\vstart\) & Prior epistemic uncertainty assigned to new competitors \\
			Maximum variance & \(\vmax\) & Upper bound on epistemic uncertainty after updates and drift \\
			Skill drift rate & \(\sdrift\) & Rate at which epistemic uncertainty grows during periods of inactivity \\
			Novelty variance & \(\psii\) & Penalty applied to topologically isolated or immature fields \\
			Tail noise variance & \(\xii\) & Score-dependent reliability penalty for deep-tail finishes (in \(W_i\) and \(R_{\mathrm{obs},i}\)) \\
			Weak-field variance & \(\omegai\) & Match-level reliability penalty for small, bottom-heavy matches \\
			Lump coefficient & \(\zetai\) & $\mathrm{Var}(P)$ in $\mathrm{s}^2$: variance of additive time lumps (PE, fumble). Not log-variance; $1/T^2$ converts it. \\
			Prediction sport variance & \(\spred\) & Sport noise used only in the prediction pipeline \\
			\bottomrule
		\end{tabular}
	\end{table}

	\begin{table}[htbp]
		\centering
		\caption{Dimensionless parameters, rates, and thresholds.}
		\label{tab:params-dimensionless}
		\begin{tabular}{llp{7cm}}
			\toprule
			\textbf{Parameter} & \textbf{Symbol} & \textbf{Role} \\
			\midrule
			Dispersion adaptation rate & \(\betad\) & Base rate for the performance dispersion exponential moving average \\
			Surprise adaptation rate & \(\gammad\) & Aggressiveness of epistemic variance inflation in response to shocks and trends \\
			Momentum adaptation rate & \(\lambdad\) & Base rate for the innovation momentum exponential moving average \\
			Pairwise blend weight & \(\alphad\) & Weight given to local pairwise opponent comparisons versus global baseline \\
			Cohort intraclass correlation & \(\rhoc\) & Assumed fraction of prior uncertainty shared among unrated competitors \\
			Student-$t$ degrees of freedom & \(\nudof\) & Heaviness of the mean-update M-estimator \\
			Student-$t$ cutoff & \(\ctcut\) & Full-strength innovation band (in sigmas) before heavy-tailed downweighting begins \\
			Baseline robustness threshold & \(\cthresh\) & Huber-style threshold (in residual standard deviations) for robust baseline estimation \\
			Tail noise start percentage & \(s_0\) & Finish percentage below which additional tail noise is applied \\
			Graph maturity threshold & \(\kmax\) & Number of events required for full topological integration into the global graph \\
			Lump reference duration & \(\tref\) & Event winning time at or above which \(\slump(T)\) is identically zero \\
			Padding floor & \(\tfloor\) & Minimum winning time triggering additive time padding \\
			Prediction behavioral scale & \(\kappap\) & Fraction of per-competitor dispersion projected into predictive uncertainty \\
			Rust decay rate & \(\lambdar\) & Strength of chronological mean reversion for inactive competitors \\
			Rust grace period & \(\tgrace\) & Inactivity (years) before mean reversion begins \\
			\bottomrule
		\end{tabular}
	\end{table}

	\subsection{Event Update}
	The following steps are the production procedure. Reliability penalties of Section~\ref{sec:robust} appear in-line.

	\subsubsection*{Step 1: Age the prior}
	Let $\Delta t$ be elapsed rating periods since the competitor's last committed update, and let $\Delta t_{\mathrm{yr}}$ be elapsed years. Epistemic drift (Section~\ref{sec:core}) and skill rust (Section~\ref{sec:rust}) are
	\begin{equation}
		V_i^{\mathrm{age}} = \min(\vmax,\; V_i + \sdrift \cdot \Delta t), \qquad
		\Delta t_{\mathrm{eff}} = \max(0,\, \Delta t_{\mathrm{yr}} - \tgrace),
	\end{equation}
	\begin{equation}
		R_i^{\mathrm{age}} = R_{\mathrm{start}} + (R_i - R_{\mathrm{start}}) \cdot \exp(-\lambdar \cdot \Delta t_{\mathrm{eff}}),
	\end{equation}
	with $R_{\mathrm{start}} = 0$. Subsequent steps use $(R_i^{\mathrm{age}}, V_i^{\mathrm{age}})$ in place of $(R_i, V_i)$.

	\subsubsection*{Step 2: Score hygiene}
	For time-plus events with $\tmin < \tfloor$, translate every time by $k = \tfloor - \tmin$ and recompute ratios from $T_i' = T_i + k$ (Section~\ref{sec:padding}). Then floor percentages and map to log scores (Section~\ref{sec:pathologies}):
	\begin{equation}
		L_i = \ln\bigl(\max(S_{\mathrm{floor}}, S_i)\bigr), \qquad S_{\mathrm{floor}} = 0.01.
	\end{equation}

	\subsubsection*{Step 3: Reliability terms and baseline weights}
	Compute the per-finish tail penalty and, for time-plus, the event-level lump penalty (Section~\ref{sec:lumps}):
	\begin{equation}
		\eta(S_i) =
		\begin{cases}
			0, & S_i \ge \szero, \\
			\xii \bigl((\szero - S_i)/\szero\bigr)^2, & S_i < \szero,
		\end{cases}
	\end{equation}
	\begin{equation}
		T_{\mathrm{eff}} = \max(\tmin,\, \epsilon), \qquad
		\slump(T) = \zetai \cdot \max\!\bigl(0,\; 1/T_{\mathrm{eff}}^2 - 1/\tref^2\bigr),
	\end{equation}
	with $\epsilon \approx 0.5$--$1.0\,\mathrm{s}$ and $\slump(T) = 0$ in hit-factor sports. Weights are
	\begin{equation}
		W_i = \frac{1}{V_i^{\mathrm{age}} + \ssport + \sigmai + \eta(S_i) + \slump(T)}.
	\end{equation}

	\subsubsection*{Step 4: Robust leave-one-out baseline}
	Residuals $a_i = R_i^{\mathrm{age}} - L_i$ estimate $\ln(\xwin)$ (Section~\ref{sec:insights}). Form $B_{\mathrm{raw}} = (\sum_i W_i a_i)/(\sum_i W_i)$ and Huber-taper (Section~\ref{sec:heavytails}):
	\begin{equation}
		r_i = \frac{a_i - B_{\mathrm{raw}}}{\sqrt{1/W_i}}, \qquad
		\tilde{W}_i = W_i \cdot \min\!\bigl(1,\, \cthresh / |r_i|\bigr).
	\end{equation}
	The leave-one-out anchor and its variance under cohort correlation $\rhoc$ (Section~\ref{sec:sparsity}) are
	\begin{equation}
		B_{-i} = \frac{\sum_{j \neq i} \tilde{W}_j a_j}{\sum_{j \neq i} \tilde{W}_j},
	\end{equation}
	\begin{equation}
		\mathrm{Var}(B_{-i}) = \frac{1}{\sum_{j \neq i} \tilde{W}_j} + \frac{\rhoc}{\bigl(\sum_{j \neq i} \tilde{W}_j\bigr)^2} \Biggl[ \Biggl( \sum_{j \neq i} \tilde{W}_j \sqrt{V_j^{\mathrm{age}}} \Biggr)^2 - \sum_{j \neq i} \tilde{W}_j^2 V_j^{\mathrm{age}} \Biggr].
	\end{equation}

	\subsubsection*{Step 5: Local blend and reconstructed performance}
	If $\alphad = 0$, set $\bobs = B_{-i}$. Otherwise select an opponent subset $\mathcal{S}_i$ (nearby finish and rating; the reference implementation also includes the top decile of the field) and form $B_{\mathrm{local}}$ and $\mathrm{Var}(B_{\mathrm{local}})$ by the same formulae as Step~4 restricted to $\mathcal{S}_i$. Sparse $\mathcal{S}_i$ should drive $\alpha_i \to 0$; otherwise $\alpha_i = \alphad$. Then
	\begin{equation}
		\bobs = (1 - \alpha_i)\, B_{-i} + \alpha_i\, \blocal, \qquad
		\mathrm{Var}(\bobs) = (1 - \alpha_i^2)\, \mathrm{Var}(B_{-i}) + \alpha_i^2\, \mathrm{Var}(\blocal),
	\end{equation}
	\begin{equation}
		O_i = L_i + \bobs.
	\end{equation}
	The cancellation $O_i \approx \ln(X_i)$ is Section~\ref{sec:core}.

	\subsubsection*{Step 6: Observation variance and quality}
	\label{step6}
	Certainty-scale dispersion, then assemble clean and total observation noise (Section~\ref{sec:gating}):
	\begin{equation}
		c_i = 1 - \min\!\bigl(1,\, V_i^{\mathrm{age}} / \vmax\bigr), \qquad
		\tilde{\sigma}_i^2 = \sigmai \cdot c_i,
	\end{equation}
	\begin{equation}
		\mu_i = \min\!\bigl(1,\, \ln(k_i+1)/\ln(\kmax+1)\bigr), \qquad
		\bar{\mu} = \frac{\sum_j \tilde{W}_j \mu_j}{\sum_j \tilde{W}_j}, \qquad
		\sigma_{\mathrm{novel}}^2 = \psii \cdot (1 - \bar{\mu}),
	\end{equation}
	where $k_i$ is the competitor's event count. Weak-field pathology (Section~\ref{sec:weakfield-overview}) is match-level.\label{sec:weakfield} With size taper $f_{\mathrm{size}}(N) = \max\bigl(0, (n_{\mathrm{max}}-N)/(n_{\mathrm{max}}-2)\bigr)$, weak non-winners $\mathcal{W} = \{j : \mathrm{place}_j > 1,\, S_j < s_w\}$, and $q = |\mathcal{W}|/(N-1)$, set $\sigma_{\mathrm{path}}^2 = 0$ if $q < q_0$ or $N \ge n_{\mathrm{max}}$; otherwise
	\begin{equation}
		g = \min\!\Biggl(1,\; \frac{1}{|\mathcal{W}| \cdot (-\ln s_w)} \sum_{j \in \mathcal{W}} \max\!\Bigl(0,\, \ln\frac{s_w}{\max(S_{\mathrm{floor}}, S_j)}\Bigr)\Biggr),
	\end{equation}
	\begin{equation}
		\sigma_{\mathrm{path}}^2 = \omegai \cdot \bigl(f_{\mathrm{size}}(N) \cdot q \cdot g\bigr).
	\end{equation}
	Then
	\begin{equation}
		R_{\mathrm{clean},i} = \ssport + \tilde{\sigma}_i^2 + \mathrm{Var}(\bobs),
	\end{equation}
	\begin{equation}
		R_{\mathrm{obs},i} = R_{\mathrm{clean},i} + \sigma_{\mathrm{novel}}^2 + \eta(S_i) + \sigma_{\mathrm{path}}^2 + \slump(T),
	\end{equation}
	\begin{equation}
		Q_i = R_{\mathrm{clean},i} / R_{\mathrm{obs},i}.
	\end{equation}

	\subsubsection*{Step 7: Trend-inflated prior}
	Using committed momentum $m_i$ from the \emph{previous} event,
	\begin{equation}
		\lambdad_i^{\mathrm{eff}} = \lambdad \cdot \max(0.25,\, c_i), \qquad
		\mathrm{Trend}_i = m_i^2 / \lambdad_i^{\mathrm{eff}},
	\end{equation}
	\begin{equation}
		V_i^{\mathrm{prior}} = \min\bigl(\vmax,\; V_i^{\mathrm{age}} + \gammad \cdot \mathrm{Trend}_i\bigr).
	\end{equation}

	\subsubsection*{Step 8: Robust Kalman mean}
	\begin{equation}
		e_i = O_i - R_i^{\mathrm{age}}, \qquad
		T_i = V_i^{\mathrm{prior}} + R_{\mathrm{obs},i}, \qquad
		K_i = V_i^{\mathrm{prior}} / T_i,
	\end{equation}
	\begin{equation}
		z_i = e_i / \sqrt{T_i}, \qquad
		w_t(z_i) = \min\!\bigl(1,\, (\nudof + \ctcut^2)/(\nudof + z_i^2)\bigr).
	\end{equation}
	The cutoff $\ctcut$ may be signed (Section~\ref{sec:heavytails}). Commit
	\begin{equation}
		R_i \leftarrow R_i^{\mathrm{age}} + K_i \cdot w_t(z_i) \cdot e_i.
	\end{equation}
	The factor $w_t$ applies only to this mean update.

	\subsubsection*{Step 9: Momentum, Shock, and posterior variance}
	\begin{equation}
		\ephys = e_i \cdot Q_i, \qquad \erobust = e_i \cdot w_t(z_i) \cdot Q_i,
	\end{equation}
	\begin{equation}
		m_i \leftarrow (1 - \lambdad_i^{\mathrm{eff}})\, m_i + \lambdad_i^{\mathrm{eff}}\, \erobust,
	\end{equation}
	\begin{equation}
		\mathrm{Shock}_i = \max\bigl(0,\; (\ephys)^2 - T_i\bigr), \qquad
		V_i \leftarrow \min\bigl(\vmax,\; V_i^{\mathrm{prior}}(1-K_i) + \gammad \cdot \mathrm{Shock}_i\bigr).
	\end{equation}
	The updated $m_i$ drives Trend on the \emph{next} event. Shock does not pass through $(1-K_i)$.

	\subsubsection*{Step 10: Dispersion EMA}
	\begin{equation}
		E_i = \max\bigl(0,\; (\ephys - m_i)^2 - V_i^{\mathrm{age}} - \ssport\bigr),
	\end{equation}
	\begin{equation}
		q_i = \mathrm{clamp}\bigl(E_i,\; 0.001\,\vstart,\; 9\cdot\max(\sigmai,\, 0.25\,\ssport)\bigr),
	\end{equation}
	\begin{equation}
		\betad_i^{\mathrm{eff}} = \betad \cdot \max(0.25,\, c_i), \qquad
		\sigmai \leftarrow \mathrm{clamp}\bigl((1-\betad_i^{\mathrm{eff}})\,\sigmai + \betad_i^{\mathrm{eff}}\, q_i,\; 0.001\,\vstart,\; \vmax\bigr).
	\end{equation}
	Use the \emph{updated} $m_i$ from Step~9 in $E_i$, so a breakout absorbed by momentum is not written into $\sigmai$.

	\section{The Prediction System and Confidence Intervals}
	\label{sec:prediction}

	Because ratings operate in log space, predicted finish ratios use log-normal bookkeeping. The system constructs a mixture distribution over plausible winners to calculate expected finish percentages and predictive confidence intervals.

	\subsection{Expected Finish Ratio}
	Let $p_w$ be the mixture weights established by the candidate probabilities. Calculating the expected finish ratio $S_i$ using a closed-form approximation requires resolving two distinct statistical boundaries: the Extreme Value behavior of the winner, and the right-truncation of the percentage scale.

	\subsubsection*{The Max Order Correction ($\Delta_w$)}
	In wide-open fields with many realistic contenders, the actual match winner will not shoot their average expected performance; they will be the maximum order statistic of the contender pool. To correct for this without requiring a full Monte Carlo simulation, we calculate the effective number of contenders using the Inverse Participation Ratio (IPR), which naturally filters out the long tail of unlikely winners:
	\begin{equation}
		N_{\mathrm{eff}} = \left(\sum_w p_w^2\right)^{-1}
	\end{equation}

	We find the expected maximum of $N_{\mathrm{eff}}$ variables using Blom's approximation. However, a naive application of Extreme Value Theory assumes that performance variance is perfectly symmetric and unbounded. In practical shooting, a competitor's performance is dominated by downside risk (mistakes, penalties, equipment malfunctions), while their maximum potential speed is biomechanically bounded by human visual processing and mechanical cycle rates.

	To preserve Extreme Value Theory assumptions regarding cohort covariance, we must strictly isolate the independent fraction of the variance. We partition the epistemic uncertainty using $(1 - \rho)$, and rely on the purely idiosyncratic prediction sport noise $\sigma_{\mathrm{pred}}^2$ to represent match-day variance. The presumed winner's rating is shifted upward by this corrected, dampened extreme-value expectation:
	\begin{equation}
		\Delta_w = \Phi^{-1}\Big(\frac{N_\mathrm{eff} - 0.375}{N_\mathrm{eff} + 0.25}\Big) \cdot \sqrt{V_w^{\mathrm{age}}(1 - \rho) + (\kappa \cdot \tilde{\sigma}_w^2) + \sigma_{\mathrm{pred}}^2}
	\end{equation}
	The adjusted pairwise log-difference between competitor $i$ and presumed winner $w$ is therefore $\tilde{\mu}_{iw} = R_i - (R_w + \Delta_w)$. For matches with a single dominant favorite ($N_{\mathrm{eff}} \approx 1$), the Blom approximation naturally evaluates to zero, preserving unbiased predictions.

	\subsubsection*{The Truncation Correction ($\Lambda_{iw}$)}
	Because $S_i$ represents a score relative to the match winner, it is strictly bounded by $S_i \le 1.0$. Conditioning the prediction on $w$ being the winner implicitly conditions the expected ratio on $L_i \le L_w$. We enforce this physical ceiling by defining a right-truncation modifier (truncated at $0$ in log space) for the log-normal expectation:
	\begin{equation}
		\Lambda_{iw} = \frac{\Phi\!\left(-\frac{\tilde{\mu}_{iw} + \sigma_{iw}^2}{\sigma_{iw}}\right)}{\Phi\!\left(-\frac{\tilde{\mu}_{iw}}{\sigma_{iw}}\right)}
	\end{equation}
	where $\Phi$ is the standard normal CDF. Because competitors share structural cohort uncertainty (formalized by the intraclass correlation $\rho$), the epistemic variance of their pairwise difference must subtract their covariance. The adjusted pairwise variance is:

	\begin{equation}
		\sigma_{iw}^2 = V_i^{\mathrm{age}} + V_w^{\mathrm{age}} - 2\rho\sqrt{V_i^{\mathrm{age}} V_w^{\mathrm{age}}} + 2\sigma_{\mathrm{pred}}^2 + \kappa(\tilde{\sigma}_i^2 + \tilde{\sigma}_w^2)
	\end{equation}

	The denominator isolates the probability mass to the universe where $w$ beats $i$, and the numerator perfectly cancels the mass of the unconstrained integral that lies above 1.0.

	\subsubsection*{The Final Expected Ratio}
	By the law of total expectation over the random winner, applying both the max order shift and the truncation boundary, the expected ratio of competitor $i$ to the winner is:
	\begin{equation}
		\hat{S}_i = (p_i \cdot 1.0) + \sum_{w \neq i} p_w \left[ \exp\!\left(\tilde{\mu}_{iw} + \frac{\sigma_{iw}^2}{2}\right) \cdot \Lambda_{iw} \right]
	\end{equation}
	The isolated $p_i$ term handles the mixture branch where $i$ is the presumed winner, yielding a ratio of exactly 1.0. For all other candidate winners, this formulation safely resolves the winner's curse bias in large fields while strictly obeying the $100\%$ scaling maximum.

	\subsection{Predictive Variance and Confidence Intervals}
	To generate predictive confidence intervals, we calculate the total log-space predictive variance using the Law of Total Variance over the mixture distribution of possible winners. The total variance is the sum of the average predictive uncertainty assuming the winner is known (Within-Component) and the variance introduced by the uncertainty of who wins (Between-Component).

	The Within-Component variance evaluates the expected pairwise uncertainty against the presumed winner. Because the predictive baseline is formed by the difference of two ratings ($R_i - R_w$), we must calculate the variance of that difference. This requires subtracting the shared cohort covariance ($\rho \sqrt{V_i V_w}$) from their combined epistemic variances.

	The sport contribution $2\sigma_{\mathrm{pred}}^2$ represents the idiosyncratic match-day noise for both competitors, and $\kappa \in [0,1]$ scales the projection of their certainty-discounted behavioral dispersions ($\tilde{\sigma}^2$). The total probability-weighted Within-Component variance is:
	\begin{equation}
		\varsigma_{i,\mathrm{within}}^2 = \sum_{w \neq i} p_w \left[ V_i^{\mathrm{age}} + V_w^{\mathrm{age}} - 2\rho\sqrt{V_i^{\mathrm{age}} V_w^{\mathrm{age}}} + 2\sigma_{\mathrm{pred}}^2 + \kappa(\tilde{\sigma}_i^2 + \tilde{\sigma}_w^2) \right]
	\end{equation}

	Because the focal competitor's expected score shifts depending on which universe occurs, we must calculate the Between-Component variance. We first establish the unconditional expected log-score (Grand Mean) across all universes:
	\begin{equation}
		\bar{L}_i = \sum_{w \neq i} p_w \tilde{\mu}_{iw}
	\end{equation}

	The Between-Component variance is the probability-weighted variance of the conditional means around this Grand Mean. In the specific universe where competitor $i$ wins, their conditional log-score is exactly $0$. The sum is therefore:
	\begin{equation}
		\varsigma_{i,\mathrm{between}}^2 = p_i (\bar{L}_i^2) + \sum_{w \neq i} p_w (\tilde{\mu}_{iw} - \bar{L}_i)^2
	\end{equation}

	The total predictive variance in log space is the sum of both components:
	\begin{equation}
		\varsigma_{i,\mathrm{total}}^2 = \varsigma_{i,\mathrm{within}}^2 + \varsigma_{i,\mathrm{between}}^2
	\end{equation}

	\subsubsection*{Moment-Matched Log-Normal Approximation}
	The predictive distribution is formally a mixture of right-truncated log-normals (one component per candidate winner $w \neq i$), plus a degenerate component at $1.0$ for the branch in which competitor $i$ is the winner. This mixture has no closed-form median or quantiles and therefore does not directly yield a multiplicative $\pm 1\sigma$ interval of known probability coverage.

	To produce a principled closed-form band, we approximate the mixture by the single log-normal $\mathrm{LN}(\mu_i^\star,\ \varsigma_{i,\mathrm{total}}^2)$ whose moments match the mixture exactly in two respects: its arithmetic mean on the ratio scale equals $\hat{S}_i$, and its variance on the log scale equals $\varsigma_{i,\mathrm{total}}^2$. Both quantities are derived above by the laws of total expectation and total variance, so no further assumptions about the mixture structure are needed. Solving the log-normal expectation identity $\hat{S}_i = \exp(\mu_i^\star + \varsigma_{i,\mathrm{total}}^2/2)$ for $\exp(\mu_i^\star)$ yields the median of the approximating distribution in ratio space:
	\begin{equation}
		\tilde{m}_i = \hat{S}_i \cdot \exp\!\left(-\frac{\varsigma_{i,\mathrm{total}}^2}{2}\right)
	\end{equation}
	The approximation is exact when the mixture collapses to a single component (e.g., when one $p_w \to 1$ and $\varsigma_{i,\mathrm{between}}^2 \to 0$) and remains tight when one favorite dominates the winner distribution. It is least faithful for wide-open fields with several equally likely winners, where the true mixture is visibly multimodal and no single log-normal can capture that shape; in such cases the approximation should be understood as an interpretable summary rather than a literal distributional claim.

	\subsubsection*{Confidence Intervals}
	Because the rating system operates natively in log space, confidence intervals in ratio space are fundamentally multiplicative. We define the geometric standard deviation of the approximating log-normal as:
	\begin{equation}
		\mathrm{GSD}_i = \exp\!\left(\sqrt{\varsigma_{i,\mathrm{total}}^2}\right)
	\end{equation}
	The $\pm 1\sigma$ predictive bounds are applied multiplicatively around the approximation median $\tilde{m}_i$:
	\begin{equation}
		S_{i,\mathrm{lower}} = \frac{\tilde{m}_i}{\mathrm{GSD}_i}, \qquad S_{i,\mathrm{upper}} = \tilde{m}_i \times \mathrm{GSD}_i
	\end{equation}
	Centering the band on $\tilde{m}_i$ rather than $\hat{S}_i$ is the standard treatment for log-normal intervals: a log-normal's $\pm 1\sigma$ probability mass is symmetric in log space around its \emph{median}, not its arithmetic mean, which sits above the median by a factor of $\exp(\varsigma_{i,\mathrm{total}}^2/2)$. Under the moment-matched approximation, the band $[\tilde{m}_i/\mathrm{GSD}_i,\ \tilde{m}_i \cdot \mathrm{GSD}_i]$ therefore carries the usual $68.27\%$ coverage interpretation; for the true mixture, the coverage is close to this value and deviates only as the mixture departs from a single log-normal.

	Notably, while both $\hat{S}_i$ and $\tilde{m}_i$ are strictly bounded to $1.0$, the upper confidence interval $S_{i,\mathrm{upper}}$ is permitted to exceed $1.0$. This is physically meaningful: an upper bound of $1.02$ indicates that a competitor's $+1\sigma$ performance would theoretically exceed the presumed winner's expected performance by 2\%, visually indicating to users that the competitor has a realistic probability of winning the match.

	\section{Comparison with Earlier Work}
	In the course of producing analysis on the shooting sports over several years, the author has read about and experimented with a number of rating algorithms.

	\subsection{Glicko-2}
	\begin{enumerate}
		\item \textbf{Native Representation of Performance Distances:} Logistic models represent the outcome of a pairwise comparison as a probability in $(0, 1)$, which is correct for binary win/loss but introduces information loss for continuous margins. A pairwise expected score of $0.95$ could correspond to a 5\% performance margin or a 50\% margin---the sigmoid saturates at the extremes, compressing large differences into a narrow band. The latent log-ratio model preserves the log-ratio structure of performance distances globally: a competitor who performs at twice the level of another has a log-gap of $\ln(2) \approx 0.69$, regardless of whether both are elite or both are novices. Pathological deep-tail updates are handled by observation-noise inflation rather than by deforming this distance scale, so prediction remains analytically simple.

		\item \textbf{Explicit Separation of Variance Sources:} Rather than routing all uncertainty through a single mechanism (e.g., Glicko-2's Rating Deviation, which conflates measurement uncertainty, skill drift, and competitor consistency), the model explicitly separates:
		\begin{itemize}
			\item $\sigma_{\mathrm{sport}}^2$: the irreducible noise of the sport for \emph{updates} (a global constant, tuned once),
			\item $\sigma_{\mathrm{pred}}^2$: idiosyncratic sport noise used only in \emph{predictions} (win probabilities and predictive bands; typically smaller than $\sigma_{\mathrm{sport}}^2$),
			\item $V_i$: measurement uncertainty in each competitor's rating (adapts via the Kalman filter),
			\item $\sigma_i^2$: per-competitor performance dispersion for \emph{updates} (certainty-scaled EMA toward a clamped target); optionally scaled by $\kappa$ for \emph{predictive} bands only,
			\item $\sigma_{drift}^2$: skill drift over time (a global constant for the aging function).
		\end{itemize}
		This decomposition makes each parameter independently interpretable and tunable. In Glicko-2, a competitor bombing a single stage can drastically inflate their RD (via the volatility $\rightarrow$ RD pathway), destabilizing future updates. Here, a single bad performance primarily affects $\sigma_i^2$ (which adapts gradually via the certainty-scaled EMA and a capped driving term) and may trigger a modest $V_i$ increase (via the innovation term), but neither mechanism produces the runaway instability that Glicko-2's volatility iteration can exhibit.
	\end{enumerate}

	\subsection{Pairwise Elo}

	\subsection{Multiplayer Elo}

	\subsection{TrueSkill and OpenSkill}

	\appendix
	\section{Parameter List}
	\label{app:parameters}

		The parameters to the Latent Log-Ratio system fall into five categories: scaling, to convert the internal log units to user-friendly display numbers; core parameters, which drive the main algorithm; parameters that improve performance on small or degenerate fields; parameters that improve observed performance when converting from ratings to predicted ratios; and chronological mean-reversion parameters for inactivity handling.

	\subsubsection*{Scaling (Display $s$, $b$)}
	\begin{itemize}
		\item \textbf{$s$, $b$ (Display Scale):} User-facing ratings use an affine map $R^{\mathrm{disp}} = s\cdot R + b$. Variance and dispersion in display space scale as $s^2$ times latent variance; display standard deviations scale as $|s|$ times latent standard deviation.
	\end{itemize}

	\subsubsection*{Core Parameters -- General}
	\begin{itemize}
		\item \textbf{$\sigma_{\mathrm{sport}}^2$ (Sport Variance):} A tuned constant representing the inherent noise of a single match observation in the sport. This is the irreducible variance that even a perfectly consistent competitor experiences due to environmental factors, range conditions, and the stochastic nature of shooting. Tuned once per sport; does not vary per competitor.
		\item \textbf{$V_{\mathrm{start}}$ (Starting Variance):} The variance assigned to competitors with no rating history, i.e. prior width for new shooters.
		\item \textbf{$V_{\max}$ (Maximum Variance):} Upper bound on committed rating variance after Kalman updates, surprise adaptation, and after drift aging.
		\item \textbf{$\sigma_{\mathrm{drift}}^2$ (Skill Drift Rate):} The rate at which rating uncertainty grows during periods of inactivity. Represents the possibility that a competitor's true skill changes when they are not being observed. Applied as $V_i \leftarrow V_i + \sigma_{\mathrm{drift}}^2 \times \Delta t$ between matches, where $\Delta t$ is measured in rating periods.
		\item \textbf{$\beta$ (Dispersion Adaptation Rate):} Base rate for the performance dispersion EMA, multiplied by a certainty factor derived from $V_i^{\mathrm{age}}/V_{\max}$ (bounded below by $0.25\,\beta$).
		\item \textbf{$\gamma$ (Surprise Adaptation Rate):} Controls how aggressively the variance $V_i$ responds to unexpectedly surprising results. Recommended range: 0.05--0.1.
		\item \textbf{$\lambda$ (Momentum Adaptation Rate):} The base rate for the innovation momentum tracker.
		\item \textbf{$\alpha$ (Pairwise Blend Weight):} Controls how much structural information from pairwise opponent comparisons is blended into the baseline-derived observed performance. When $\alpha = 0$, the model reduces to the pure baseline model.
		\item \textbf{$\nudof$ (Student-$t$ Degrees of Freedom):} Heaviness of the mean-update M-estimator \(w_t(z)\). Smaller \(\nudof\) damps moderate outliers more aggressively; \(\nudof \to \infty\) recovers an undamped Gaussian mean update. Default: \(4\).
		\item \textbf{$\ctcut$ (Student-$t$ Cutoff):} Full-strength innovation band, in residual sigmas, before \(w_t\) begins to taper. May be split by the sign of the innovation. Setting \(\ctcut = 1\) recovers the standard Student-$t$ kernel.
	\end{itemize}

	\subsubsection*{Core Parameters -- Off-Diagonal Covariance Approximation}
	\begin{itemize}
		\item \textbf{$\rho$ (Cohort Intraclass Correlation):} A global correlation coefficient representing the fraction of epistemic uncertainty assumed to be shared structurally among competitors (e.g., a local club whose ratings are accurate relative to each other, but floating relative to the global mean). Because this covariance mathematically scales with the absolute epistemic variance of the field, it exerts massive influence on new cohorts but gracefully vanishes for established veterans. In practice, $\rho$ prevents certainty collapse (the ``Day 1'' rating explosion) by explicitly bounding the baseline variance to a cohort asymptote, preventing the filter from falsely trusting an uncertain field's average as a rigid mathematical datum.
		\item \textbf{$k_{\mathrm{max}}$ (Graph Maturity Threshold):} The number of historical events at which a competitor is assumed to be fully integrated into the global network.
		\item \textbf{$\psi^2$ (Novelty Variance):} The maximum noise penalty applied to topologically isolated, immature fields. Along with $k_{\mathrm{max}}$, it limits the impact of extreme performance margins shot in low-history networks: ``The margins between these competitors are directionally meaningful, but their absolute magnitudes are untrustworthy as a measure of global top-end skill.'' Fundamentally, $k_{\mathrm{max}}$ and $\psi^2$ prevent the filter from falsely trusting the field's structural variance.
	\end{itemize}

	\subsubsection*{Small and Degenerate Fields}
	\begin{itemize}
		\item \textbf{$c$ (Baseline Robustness Threshold):} A Huber-style threshold, measured in residual standard deviations, used to downweight extreme field anchors when estimating match baseline.
		\item \textbf{$s_0$ (Tail Noise Start Percentage):} The finish percentage below which very weak results begin receiving additional observation noise. Above $s_0$, the ordinary variance model is used.
		\item \textbf{$\xi^2$ (Tail Noise Variance):} The maximum additional observation variance applied to the deepest tail. Larger values make the model more skeptical of very weak finishes without changing the log-ratio mean link.
		\item \textbf{$\omega^2$ (Weak-Field Variance):} The maximum additional match-context observation variance applied when a field is both small and bottom-heavy.
		\item \textbf{$n_{\mathrm{max}}$ (Weak-Field Max Size):} The field size at or above which weak-field damping decays to zero.
		\item \textbf{$s_w$ (Weak Finish Threshold):} The finish threshold below which a non-winning score counts as ``weak'' for weak-field detection.
		\item \textbf{$q_0$ (Weak Fraction Threshold):} The minimum fraction of weak non-winners required before the weak-field term activates.
		\item \textbf{$\zetai$ (Lump Coefficient):} $\mathrm{Var}(P)$ in clock seconds squared: the variance of additive time lumps (a procedural, a fumbled reload, an isolated miss), \emph{not} $\mathrm{Var}(P/T)$. Combined with $\tref$, $\slump(T) = \zetai \cdot \max(0, 1/T_{\mathrm{eff}}^2 - 1/\tref^2)$ is the excess log-variance of those lumps on a short event relative to a medium course. Set $\zeta$ to the RMS blunder in seconds and square it (a $2\,\mathrm{s}$ PE as a typical lump $\Rightarrow$ $\zetai \approx 4\,\mathrm{s}^2$; use a smaller value if such lumps are rare).
		\item \textbf{$\tref$ (Lump Reference Duration):} Event winning time at or above which $\slump(T)$ is identically zero (typically $25$--$30\,\mathrm{s}$).
		\item \textbf{$\tfloor$ (Padding Floor):} Minimum winning time; below it, additive translation restores positivity before logarithms are taken.
	\end{itemize}

	\subsubsection*{Prediction Parameters}
	\begin{itemize}
		\item \textbf{$\sigma_{\mathrm{pred}}^2$ (Prediction Sport Variance):} Idiosyncratic sport noise used only in the prediction pipeline (win-probability denominators and predictive variance).
		\item \textbf{$\kappa$ (Prediction Behavioral Scale):} Dimensionless coefficient in $[0,1]$ that scales how much per-competitor behavioral variance $\sigma_i^2$ enters \emph{posterior predictive} variance for finish-ratio bands.
	\end{itemize}

	\subsubsection*{Mean Reversion Parameters}
	\begin{itemize}
		\item \textbf{$t_{\mathrm{grace}}$ (Rust Grace Period):} Inactivity duration, measured in years, before chronological mean-reversion begins. For elapsed time $\Delta t$ since last match, decay uses $\Delta t_{\mathrm{eff}} = \max(0, \Delta t - t_{\mathrm{grace}})$.
		\item \textbf{$\lambda_{\mathrm{rust}}$ (Rust Decay Rate):} Exponential decay coefficient, measured in inverse years, controlling how quickly inactive competitors regress toward baseline after grace time. Applied through $\exp(-\lambda_{\mathrm{rust}} \cdot \Delta t_{\mathrm{eff}})$ as described in Section~\ref{sec:rust}.
	\end{itemize}

	\section{Practical Initialization and Tuning Guide}
	\label{app:tuning}

	Because the Latent Log-Ratio model decomposes uncertainty into strictly separated physical mechanisms, attempting to tune the model intuitively using human-readable standard deviations often leads to severe pathologies (e.g., the ``Square Root Deception,'' where users misjudge the gravitational pull of precision-weighted means).

	Tuning must follow a structured, normalized sequence. The irreducible match noise of the sport acts as the fundamental physical anchor ($1.0\,\mathrm{SV}$). All epistemic and aleatoric variances are set as multiples of this unit, followed by adjusting the dimensionless adaptation rates.

	\subsection{Phase 1: Determine the Anchor ($1.0\,\mathrm{SV}$)}
	The Sport Variance ($\sigma_{\mathrm{sport}}^2$) represents the inherent, irreducible background noise of a single match observation in the sport (environmental factors, lighting, stage design).
	\begin{enumerate}
		\item Identify a cohort of highly consistent, active veterans.
		\item Measure the steady-state variance of their log-score residuals.
		\item Base $\sigma_{\mathrm{sport}}^2 \equiv 1.0\,\mathrm{SV}$ on this empirical variance. If a consistent pro fluctuates by $\pm 3\%$ match-to-match, $\sqrt{1.0\,\mathrm{SV}} \approx 0.03$.
	\end{enumerate}

	Because this empirical measurement inherently contains its own estimation error, and indeed is challenging to measure in sports where environmental factors like weather or course difficulty disguise the underlying signal, small adjustments to $\sigma_{\mathrm{sport}}^2$ may be practically useful. By tweaking \emph{just} the Sport Variance (while holding the epistemic parameters like $V_{\mathrm{start}}$ and $\sigma_{\mathrm{drift}}^2$ at their established absolute values), it acts as a thermodynamic ``temperature'' dial for the model:
	\begin{itemize}
		\item \textbf{Raising Sport Variance (Cooling):} The model assumes the environment is noisier than strictly measured. Total expected uncertainty ($T_i$) grows relative to the immutable physical match residuals ($e_i$). This drops the Kalman gain slightly and makes the non-linear Variance Surprise ($e_i^2 > T_i$) harder to trigger, resulting in a smoother, less reactive filter. Moving too far in this direction will yield a filter that fails to update on true changes to the underlying competitor variables.
		\item \textbf{Lowering Sport Variance (Heating):} The model assumes a more sterile environment than measured. $T_i$ shrinks, making the filter more reactive. Moving too far in this direction will yield highly volatile ratings, over-updating on normal fluctuations.
	\end{itemize}

	\subsection{Phase 2: Establish the Golden Ratios (Units of $\mathrm{SV}$)}
	With the baseline $1.0\,\mathrm{SV}$ defined, initialize the system's variance-space parameters as relative multiples of this anchor.

	\begin{itemize}
		\item \textbf{Skill Drift Rate ($\sigma_{\mathrm{drift}}^2$): $0.5\,\mathrm{SV}$ per year (Range: $0.1 - 1.0$).}
		Adult motor skills degrade slowly. Half a match's worth of uncertainty per year allows veterans to remain stable while preventing absolute certainty collapse.
		\item \textbf{Starting Variance ($V_{\mathrm{start}}$): $25.0\,\mathrm{SV}$ (Range: $20.0 - 40.0$).}
		New competitors require enough variance to traverse the demographic distribution rapidly. If competitor graphs often show an initial spike followed by decay, tune down slightly.
		\item \textbf{Maximum Variance ($V_{\max}$): $33.0\,\mathrm{SV}$ (Range: $30.0 - 50.0$).}
		An absolute computational ceiling slightly higher than $V_{\mathrm{start}}$, allowing returning shooters or massive variance surprises to snap the Kalman gain back to near 1.0 without overflowing the arithmetic. Note that the certainty parameter $c_i$ defined in Step~6 of Section~\ref{sec:algorithm} scales the update speed of dispersion $\sigma_i^2$ and momentum $m_i$, so $V_{\max}$ should be moved in proportion with $V_\mathrm{start}$ to prevent unintentional changes to the movement of other parameters.
		\item \textbf{Initial Dispersion ($\sigma_{i,\mathrm{init}}^2$): $0.1\,\mathrm{SV}$ (Range: $0.05 - 0.25$).}
		Competitors must be treated as consistent until proven otherwise. Initializing dispersion at a fraction of $\mathrm{SV}$ ensures the Kalman denominator remains small, allowing unhindered early rating convergence.
		\item \textbf{Tail Noise ($\xi^2$): $\approx 33.0\,\mathrm{SV}$:}
		Requires empirical tuning based on the sport's deep-tail geometry. Set high enough to collapse the update weight of catastrophic failures to near-zero, if needed.
		\item \textbf{Weak-Field Variance ($\omega^2$): $\approx 200.0\,\mathrm{SV}$.}
		A massive structural penalty applied to degenerate, small-and-weak fields, ensuring they provide virtually no mathematically rigid evidence about elite top-end skill separation.
		\item \textbf{Novelty Variance ($\psi^2$): $5.0 \mathrm{SV}$ (Range: $1.0 - 25.0$).}
		Novelty variance serves as a light regularizing factor on fields of predominantly new competitors. It targets high-level competitors who appear and retire early in a dataset, earning incorrectly high ratings against a competitor graph that is not yet widely connected. $5.0$ represents minimal damping, but is experimentally sufficient to correct most unexpected outliers without unduly slowing the number of events required for new groups of ratings to converge.
		\item \textbf{Prediction Sport Variance ($\sigma_{\mathrm{pred}}^2$): $0.0\,\mathrm{SV}$.}
		Philosophically, predictive uncertainty should scale with the competitors' actual historical behaviors. Set this to zero and use $\kappa$ (see Phase 3) to widen predictive error bars. Only inject flat $\sigma_{\mathrm{pred}}^2$ if $\kappa$ is fully saturated at $1.0$ and predictions remain too narrow, or if the underlying conditions of the sport produce idiosyncratic variance that cannot be attributed to differences in competitor skill.
	\end{itemize}

	\subsection{Phase 3: Dimensionless Parameters and Thresholds}
	Finally, tune the adaptation rates, blend weights, and Huber thresholds.

	\begin{itemize}
		\item \textbf{Prediction Behavioral Scale ($\kappa$):} \textbf{(Range: $0.0 - 1.0$).} Use as the primary tool for predictive error bar size control. It scales how much of a competitor's behavioral dispersion ($\sigma_i^2$) is projected into their published predictive bands.
		\item \textbf{Cohort Intraclass Correlation ($\rho$): $0.3$ (Range: $\ge 0$).}
		Prefer the lowest value that reliably eliminates the ``Day 1 Rating Explosion'' (the artifact where unrated fields collapse baseline variance to zero). The observed symptom of this problem is early competitors in weak or unknown fields climbing to the top of the ratings, while empirically higher-skilled competitors struggle to reach the same region after the cold start has completed. Increasing this parameter will damp early changes in any new group of ratings by increasing observational noise and thereby decreasing the Kalman gain.
		\item \textbf{Momentum Adaptation Rate ($\lambda$): $0.2$}
		Controls the EMA for tracking directional skill trends. The half-life of an exponential moving average is given by $ln(0.5) / ln(1 - \alpha)$.
		The default yields a half-life of $\approx 3.1$ events. A relatively short half-life (according to the schedule of the sport) yields faster reaction to both the beginning and end of a streak.
		\item \textbf{Dispersion Adaptation Rate ($\beta$): $0.1$.}
		Controls the EMA for behavioral volatility.The default value of $0.1$ yields a half-life of $\approx 6.6$ events. A long window relative to the sport's schedule (and to $\lambda$) prevents too-eager increases in dispersion, which can lead to the ``aleatoric armor'' pathology.
		\item \textbf{Surprise Adaptation Rate ($\gamma$): $0.1$.}
		Controls how aggressively $V_i$ inflates when the momentum or squared error breaches expected variance. Depends on event frequency and the nature of skill shocks in the sport.
		\item \textbf{Pairwise Blend Weight ($\alpha$): $0.1$ (Range: $0.0 - 1.0$).}
		Empirically, precision-weighted global baselines are more reliable than sparse pairwise links. However, blending a small amount of pairwise data ($\alpha=0.1$) allows the system to catch localized environmental signals (e.g., competitors scheduled for different days with different weather) or asymmetries in field composition (e.g., a weak center but a strong top end) that the global baseline smooths over.
		\item \textbf{Student-$t$ Degrees of Freedom ($\nudof$): $4$ (Range: $2$--$10$).}
		Controls how fast \(w_t\) falls once \(|z|\) exceeds \(\ctcut\). \(\nudof = 4\) is the usual robust-$t$ default (\(\approx 29\%\) damping at \(z = -1.75\) when \(\ctcut = 1\)). Raising \(\nudof\) toward \(8\)--$10$ softens the taper.
		\item \textbf{Student-$t$ Cutoff ($\ctcut$): $2$ (Range: $0$--$3$; optionally signed).}
		Where the taper \emph{starts}. \(1.0\) begins immediately past \(1\sigma\); \(2.0\) lets typical matches through; a smaller downside cutoff (e.g.\ \(1.0\)) with a larger upside cutoff is the one-sided extension when contamination is known to be negative. \(\ctcut\) and \(\nudof\) interact: a tight cutoff with large \(\nudof\) still barely damps.
		\item \textbf{Baseline Robustness ($c$): $2.5$.}
		A Huber-style threshold marking the boundary between $L_2$ (squared) and $L_1$ (absolute) penalty weighting. A value of 2.5 means the filter trusts residuals normally up to 2.5 standard deviations, but linearly downweights the influence of outliers beyond that. Lower values (e.g., 2.0) aggressively clip outliers in highly volatile sports; higher values (3.0+) trust the field's arithmetic mean more strictly.
		\item \textbf{Tail Noise Start ($s_0$): $0.6$.}
		The finish percentage below which the $\xi^2$ tail noise begins applying. Tune depending on the severity of the winner's curse and the geometry of the sport's deep tail.
		\item \textbf{Weak-Field Thresholds ($n_{\mathrm{max}}, s_w, q_0$):}
		Highly dependent on sport demographics. Define the size and shape of a match that should be legally penalized by $\omega^2$.
		\item \textbf{Graph Maturity Threshold ($k_\mathrm{max}$): $8.0$.}
		How quickly novelty variance $\psi^2$ decays as precision-weighted field history length grows. Dependent on sport demographics. For values between 5 and 20, the penalty reaches half strength at approximately $0.2 \cdot k_\mathrm{max}$.
		\item \textbf{Lump Coefficient ($\zetai$) and Reference Duration ($\tref$):}
		Time-plus only. $\zetai$ is $\mathrm{Var}(P)$ in $\mathrm{s}^2$, not a log-variance. Set $\zeta$ to the RMS additive blunder you want treated as geometry (a $2\,\mathrm{s}$ PE $\Rightarrow$ start near $4\,\mathrm{s}^2$; lower if those events are rare). Set $\tref$ to a medium course where a lump of that size is already diluted ($20$--$30\,\mathrm{s}$). Check that a $2\,\mathrm{s}$ gap on an $8\,\mathrm{s}$ stage is a few observation SDs of $R_{\mathrm{obs}}$, not $7\sigma$ of $\ssport$ alone, and that $\slump(T)=0$ on ordinary field courses.
	\end{itemize}

	\bibliographystyle{ieeetr}
	\bibliography{paper}
\end{document}